\section{Framework Evolution}
\label{appendix:evolution}

This framework did not evolve through theoretical speculation; it hardened through adversarial engagement with live systems. Each version update represents the closure of a specific "loophole" used by operators to claim public legitimacy while retaining private control.

\begin{description}[style=nextline, leftmargin=0pt, labelindent=0pt]

\item[\textbf{Version 1.0 (Definition)}]
Established the baseline thesis: Corrigibility is a structural precondition, distinct from feature sets. It defined the five tests to answer the foundational question: \textit{What distinguishes public infrastructure from private monopoly?}

\item[\textbf{Version 1.2 (Assessment)}]
\textbf{Closed the Self-Attestation Loophole.} Early trials showed operators self-declaring success. This version separated the \textit{Infrastructure Manifest} from the \textit{Corrigibility Assessment}, establishing the adversarial verification model where publicness must be proven by independent parties.

\item[\textbf{Version 1.4 (Logic)}]
\textbf{Closed the Ambiguity Loophole.} Eliminated partial credit and "maturity models." The shift to strictly binary logic established that failure in any single dimension breaks the feedback loop required by Ashby's Law, rendering the system unstable.

\item[\textbf{Version 1.5 (Invariance)}]
\textbf{Closed the Exceptionalism Loophole.} Extended the framework to AI and decentralized ledgers, demonstrating that neither statistical opacity nor cryptographic immutability exempts a system from the requirements of Reversibility (EXIT) and Correction (GOVERN).

\item[\textbf{Version 1.6 (Constraint)}]
\textbf{Closed the Consultation Loophole.} Redefined GOVERN as non-overrideable internal constraint rather than advice or redress. Explicitly disqualified post-execution remedies and advisory bodies from constituting system governance.

\item[\textbf{Version 1.6.2 (Reproduction)}]
\textbf{Closed the Substitution Loophole.} Distinguished EXIT (market choice) from FORK (reproduction). Switching landlords is not forking infrastructure.

\item[\textbf{Version 1.7 (Layering)}]
\textbf{Closed the DID-Washing Loophole.} Added layer-decomposed evaluation requiring analysis at each architectural stratum. Systems claiming corrigibility through identifier decentralization while retaining centralized governance at higher layers are now explicitly diagnosed. Introduced the propagation rule: $T_{\text{system}} = \bigwedge T(L_i)$.

\end{description}

% ============================================
% REFERENCES (at end, per academic standard)
% ============================================
\bibliographystyle{plainnat}
\begin{thebibliography}{99}

\bibitem[Ashby(1956)]{ashby1956}
Ashby, W. Ross.
\newblock {\em An Introduction to Cybernetics}.
\newblock Chapman \& Hall, 1956.

\bibitem[Ostrom(1990)]{ostrom1990}
Ostrom, Elinor.
\newblock {\em Governing the Commons: The Evolution of Institutions for Collective Action}.
\newblock Cambridge University Press, 1990.

\bibitem[Stallman(2002)]{stallman2002}
Stallman, Richard.
\newblock {\em Free Software, Free Society: Selected Essays}.
\newblock GNU Press, 2002.

\bibitem[Winner(1980)]{winner1980}
Winner, Langdon.
\newblock Do Artifacts Have Politics?
\newblock {\em Daedalus}, 109(1):121--136, 1980.

\bibitem[Hirschman(1970)]{hirschman1970}
Hirschman, Albert O.
\newblock {\em Exit, Voice, and Loyalty: Responses to Decline in Firms, Organizations, and States}.
\newblock Harvard University Press, 1970.

\bibitem[Lessig(2006)]{lessig2006}
Lessig, Lawrence.
\newblock {\em Code: Version 2.0}.
\newblock Basic Books, 2006.

\bibitem[Bowker and Star(1999)]{bowker1999}
Bowker, Geoffrey C. and Star, Susan Leigh.
\newblock {\em Sorting Things Out: Classification and Its Consequences}.
\newblock MIT Press, 1999.

\bibitem[G20(2023)]{g20}
G20.
\newblock New Delhi Leaders' Declaration.
\newblock G20 Summit, September 2023.

\bibitem[UNDP(2024)]{undp}
United Nations Development Programme.
\newblock Digital Public Infrastructure: Definitions and Core Concepts.
\newblock UNDP Digital Strategy Office, 2024.

\bibitem[Gebru et al.(2021)]{gebru2021}
Gebru, Timnit, Morgenstern, Jamie, Vecchione, Briana, et al.
\newblock Datasheets for Datasets.
\newblock {\em Communications of the ACM}, 64(12):86--92, 2021.

\bibitem[Mitchell et al.(2019)]{mitchell2019}
Mitchell, Margaret, Wu, Simone, Zaldivar, Andrew, et al.
\newblock Model Cards for Model Reporting.
\newblock In {\em Proceedings of the Conference on Fairness, Accountability, and Transparency}, pages 220--229, 2019.

\bibitem[Bhatia(2019)]{bhatia2019}
Bhatia, Gautam.
\newblock {\em The Transformative Constitution: A Radical Biography in Nine Acts}.
\newblock HarperCollins India, 2019.

\bibitem[Leveson(2011)]{leveson2011}
Leveson, Nancy G.
\newblock {\em Engineering a Safer World: Systems Thinking Applied to Safety}.
\newblock MIT Press, 2011.

\bibitem[Khera(2019)]{khera2019}
Khera, Reetika.
\newblock {\em Dissent on Aadhaar: Big Data Meets Big Brother}.
\newblock Orient BlackSwan, 2019.

\bibitem[Cohen(2019)]{cohen2019}
Cohen, Julie E.
\newblock {\em Between Truth and Power: The Legal Constructions of Informational Capitalism}.
\newblock Oxford University Press, 2019.

\bibitem[Supreme Court of India(2017)]{puttaswamy2017}
Supreme Court of India.
\newblock Justice K.S. Puttaswamy (Retd.) v. Union of India.
\newblock Writ Petition (Civil) No. 494 of 2012, (2017) 10 SCC 1.

\bibitem[Collingridge(1980)]{collingridge1980}
Collingridge, David.
\newblock {\em The Social Control of Technology}.
\newblock Frances Pinter, 1980.

\bibitem[Pasquale(2015)]{pasquale2015}
Pasquale, Frank.
\newblock {\em The Black Box Society: The Secret Algorithms That Control Money and Information}.
\newblock Harvard University Press, 2015.

\bibitem[W3C(2022)]{w3cdid2022}
World Wide Web Consortium.
\newblock Decentralized Identifiers (DIDs) v1.0.
\newblock W3C Recommendation, July 19, 2022.
\newblock \url{https://www.w3.org/TR/did-core/}

\bibitem[European Union(2024)]{eidas2024}
European Union.
\newblock Regulation (EU) 2024/1183 amending Regulation (EU) No 910/2014 (eIDAS 2.0).
\newblock {\em Official Journal of the European Union}, 2024.

\bibitem[Khera(2017)]{khera2017}
Khera, Reetika.
\newblock Impact of Aadhaar on Welfare Programmes.
\newblock {\em Economic and Political Weekly}, 52(50):61--70, 2017.

\bibitem[European Commission(2021)]{eu2021oss}
European Commission.
\newblock Study about the Impact of Open Source Software and Hardware on Technological Independence, Competitiveness and Innovation in the EU Economy.
\newblock European Commission, Directorate-General for Communications Networks, Content and Technology, 2021.

\end{thebibliography}